% ─────────────────────────────────────────────────────────────────────────────
\section{Final Project Status, Limitations, and Future Work}
\label{sec:status}
% ─────────────────────────────────────────────────────────────────────────────

The project concluded with a working implementation at prototype level. The core
hybrid pipeline was completed conceptually and functionally: detection, region isolation,
reasoning, structured output, and closed-network design were all part of the final system.
What remains unfinished is not the basic architecture, but the broader packaging and
qualification work required for a fully released or field-hardened product.

\subsection{Final Status Overview}

\begin{table}[H]
  \centering
  \caption{Final project status by work package}
  \begin{tabularx}{\textwidth}{@{}p{5.2cm}p{2.2cm}X@{}}
    \toprule
    Work package & Status & Notes \\
    \midrule
    Problem analysis and literature review & Complete & Background study established the
    need for a detector-plus-reasoner architecture. \\
    Requirements and constraints definition & Complete & Functional, non-functional, and
    ethical constraints were documented early and preserved in the final scope. \\
    Two-stage architecture design & Complete & The final system architecture was fixed as
    ``Gözcü + Analist.'' \\
    Working prototype implementation & Complete & Core detection, cropping, reasoning,
    and structured output workflow reached prototype form. \\
    Closed-network deployment concept & Complete & The runtime design avoids external
    cloud dependence and supports local execution. \\
    Archived automated benchmark evidence & Pending & Final report relies on retained
    system-level validation rather than a preserved benchmark bundle. \\
    \bottomrule
  \end{tabularx}
\end{table}

\subsection{Main Limitations}

The final prototype has several important limitations.

\begin{itemize}[nosep]
  \item \textbf{Fine-grained thermal ambiguity:} exact make/model recognition and highly
    specific semantic claims remain unreliable when the crop itself contains limited detail.
  \item \textbf{Dataset dependence:} output quality depends heavily on the representativeness
    of the thermal training and adaptation data.
  \item \textbf{Edge-compute pressure:} a strong VLM is useful but computationally expensive,
    especially when low latency is required.
  \item \textbf{Precision-performance trade-off:} FP16, quantized, or engine-compiled
    execution paths can improve throughput, but they must be checked carefully against any
    loss of semantic fidelity or detector stability.
  \item \textbf{Packaging complexity:} cross-architecture builds, container runtime mapping,
    and target-specific dependency management increase engineering effort even when the
    high-level model logic is already working.
\end{itemize}

These limitations do not invalidate the project, but they define its maturity level accurately.

\subsection{Low-Level Engineering Constraints}

The low-level deployment view introduces several practical constraints that are less visible in
high-level architecture diagrams.

\begin{itemize}[nosep]
  \item Increasing semantic depth usually increases latency and RAM pressure.
  \item More aggressive visual enrichment strategies, such as multiple crops or repeated target
    passes, may improve interpretive quality but are expensive on smaller embedded devices.
  \item Detector acceleration and reasoning acceleration do not scale identically; a pipeline may
    have a fast first stage and still bottleneck on the second stage.
  \item Visualization and local container execution can require explicit host-device integration,
    which is operationally simple in principle but non-trivial to standardize robustly.
\end{itemize}

These are implementation realities rather than conceptual flaws, but they directly affect what
counts as a successful edge deployment.

\subsection{Future Work}

The most useful future extensions are:

\begin{enumerate}[nosep]
  \item Build a richer thermal instruction-tuning dataset for the reasoning stage, especially
  around vehicle subtype, weapon clue, and posture interpretation.
  \item Add stricter structured prompting and schema validation to reduce overly broad or
  weakly grounded textual interpretations.
  \item Perform repeatable latency, memory, and stability measurements on the final target edge
  hardware under longer-running workloads.
  \item Publish a reproducible containerized deployment workflow, including Jetson-targeted
  runtime configuration, startup instructions, and optimization settings.
  \item Extend the output layer toward mission logging, audit trails, and more formal user-
  interface integration.
\end{enumerate}

\subsection{Overall Final Assessment}

Despite the limitations above, the final assessment of the project is positive. The project
successfully demonstrated that a hybrid architecture is a practical way to move from thermal
detection toward thermal understanding in a closed-network environment. That architectural
result is the main contribution of the work.
